%!TEX root = ../../main.tex
% ============================================================
% 数学\几何作图\六年级.tex
% 本文件由 main.tex 通过 \input 引入，请勿单独编译
% ============================================================


\section{解答：网格图的图形变换（平移、旋转与轴对称）}

\vspace{0.4cm}

\noindent\textbf{2. (1) 把三角形向左平移 5 格。}\\
\noindent\textbf{   (2) 把长方形绕点 $O$ 逆时针旋转 $90^{\circ}$。}\\
\noindent\textbf{   (3) 把右边的图形补全,使它成为轴对称图形。}

\vspace{1.2cm}

\begin{center}
% =====================================================================
% 【整体绘制思路：网格图变换综合应用】
% 1. 基础网格：为了保证比例一致性，绘制一个 31x8 的虚线网格系统。外围用粗灰线包围。
% 2. 坐标映射原图(纯黑色)：
%    - 长方形：左下角点 O 在 (6,2)，右上角 (9,6)。
%    - 三角形：三个顶点分别为 左下(16,3)，上(18,6)，右下(19,2)。
%    - 轴对称基础图形(半个飞镖)：(26,6) - (25,4) - (22,3) - (25,2) - (26,0)。
%    - 对称轴：位于 x=26 的点划线。
% 3. 绘制变化后的答案图形(加粗红色 red!70!black)：
%    - 三角形平移：所有横坐标 -5 -> (11,3), (13,6), (14,2)。
%    - 长方形旋转：绕(6,2)逆时针旋转 90 度 -> 生成一个向左倒的新长方形 (2,2) 到 (6,5)。
%    - 轴对称补全：原图向右沿 x=26 进行镜像映射 -> (26,6) - (27,4) - (30,3) - (27,2) - (26,0)。
% =====================================================================
\begin{tikzpicture}[>=Stealth, scale=0.45]
  % 1. 绘制网格线
  % 内部虚线
  \draw[help lines, dashed, step=1] (0, 0) grid (31, 8);
  % 外部实线边框
  \draw[thick, gray!80!black] (0, 0) rectangle (31, 8);
  
  % 2. 绘制已知原始条件 (黑色线条)
  % 长方形及点O标记
  \draw[thick, black] (6, 2) rectangle (9, 6);
  \node[anchor=north east, inner sep=2pt, font=\small] at (6, 2) {$O$};
  
  % 三角形
  \draw[thick, black] (15, 5) -- (18, 8) -- (19, 4) -- cycle;
  
  % 右侧半个轴对称图形 (左翼)
  \draw[thick, black] (26, 6) -- (25, 4) -- (22, 3) -- (25, 2) -- (26, 0);
  % 对称轴
  \draw[dash dot, thick, black] (26, -0.3) -- (26, 8.3);
  
  % 3. 绘制题目要求的作答结果 (深红色加粗突出显示)
  % (1) 三角形向左平移 5 格 (x坐标全部减 5)
  \draw[thick, red!70!black, line width=1.2pt] (10, 5) -- (13, 8) -- (14, 4) -- cycle;
  
  % (2) 长方形绕点 O 逆时针旋转 90 度
  % 原始宽度 3, 高度 4。绕左下角旋转90度后，变为宽度4，高度3，且位于该点左侧。
  \draw[thick, red!70!black, line width=1.2pt] (2, 2) rectangle (6, 5);
  
  % (3) 补全成轴对称图形 (基于 x=26 对称)
  % 对应点坐标 x' = 26 + (26 - x) => x' = 52 - x
  % (25,4)->(27,4), (22,3)->(30,3), (25,2)->(27,2)
  \draw[thick, red!70!black, line width=1.2pt] (26, 6) -- (27, 4) -- (30, 3) -- (27, 2) -- (26, 0);

\end{tikzpicture}
\end{center}

\vspace{0.8cm}

{\small\color{gray}
\textbf{解析：}\\
1. \textbf{平移}：确定平移方向（向左）和距离（5格）。找出原图形的几个关键顶点，分别把这些点向左平移5格后，依次连接起来即可。\\
2. \textbf{旋转}：确定旋转中心点（点 $O$），旋转方向（逆时针），旋转角度（$90^{\circ}$）。原长方形位于 $O$ 点右上方，底边长3格，竖边长4格。绕点 $O$ 逆时针旋转 $90^{\circ}$ 后，原先水平向右的底边变成竖直向上（长3格），原先竖向的边变成水平向左（长4格），形成一个位于点 $O$ 左上方、长为4、宽为3的新长方形。\\
3. \textbf{轴对称}：找准对折的“脊梁”（中心的点划线作为对称轴）。找出左边半个图案的三个拐点，依据“对称点到对称轴的距离相等”这一特点，在虚线右侧找出对应的三个落脚点，用线段顺次连接。
}

\vspace{0.6cm}

{\small\color{blue!70!black}
\textbf{绘图基本原理与步骤：}\\
\textbf{1. 网格化与坐标系建立}：在 TikZ 中利用 \texttt{grid} 命令生成一个宽 31 单位、高 8 单位的虚线网格骨架（底角设定为坐标原点 \texttt{(0,0)}），方便后续所有几何图形全部依据整数网络点吸附作图。\\
\textbf{2. 坐标域数学变换法则}：所有结果图形均由严格坐标运算得出：
   \begin{itemize}
       \item \textbf{平移映射}：将原三角形节点集 $\{(15,5), (18,8), (19,4)\}$ 按向量 $(-5, 0)$ 叠加，生成平移多边形。
       \item \textbf{旋转方程}：原矩阵矩形宽 3 高 4 的向量簇绕基点 $O=(6,2)$ 逆时针旋转 $90^{\circ}$，精确生成左置排布的新矩形域坐标。
       \item \textbf{镜像反射}：基于对称轴 $x=26$ 进行轴对称算子变换（公式 $x' = 2 \times 26 - x$），精密计算出对应极值点 $(27,4), (30,3), (27,2)$ 完成右侧飞镖图补全。
   \end{itemize}
\textbf{3. 作答视觉区分}：已知题干条件维持 \texttt{black} 常规粗线，变化产生的结果则采用辅导教材专属高亮答案红（\texttt{red!70!black}，搭配 \texttt{1.2pt} 线宽）进行叠加式渲染图层设计。
}


\newpage



\section{解答：按比例放大或缩小图形}

\vspace{0.4cm}

\noindent\textbf{3. (1) 按 1：3 的比画出正方形缩小后的图形。}\\
\noindent\textbf{   (2) 按 2：1 的比画出三角形放大后的图形。}

\vspace{1.2cm}

\begin{center}
% =====================================================================
% 【整体绘制思路：比例缩放作图】
% 1. 基础网格：根据长宽比预估，创建一个宽31, 高10的网格空间。
% 2. 原始黑色图形：
%    - 左侧大正方形：边长 6 格。设定坐标 (2,2) 到 (8,8)。
%    - 右侧小直角三角形：竖直高 3 格，水平宽 2 格。设定坐标 (16,5)-(16,8)-(18,8)。
% 3. 目标红色图形 (red!70!black, line width=1.2pt)：
%    - 正方形按 1:3 缩小：新边长为 6/3 = 2。紧贴在大正方形右侧，顶端对齐。坐标 (10,6) 到 (12,8)。
%    - 三角形按 2:1 放大：新高 3*2 = 6，新宽 2*2 = 4。紧贴在小三角形右侧，顶端对齐。坐标 (20,2)-(20,8)-(24,8)。
% =====================================================================
\begin{tikzpicture}[>=Stealth, scale=0.45]
  % 1. 绘制网格线
  \draw[help lines, dashed, step=1] (0, 0) grid (31, 10);
  \draw[thick, gray!80!black] (0, 0) rectangle (31, 10);
  
  % 2. 绘制原图 (黑色粗线)
  % 大正方形 (6x6)
  \draw[thick, black] (2, 2) rectangle (8, 8);
  % 小直角三角形 (高3, 宽2)
  \draw[thick, black] (16, 5) -- (16, 8) -- (18, 8) -- cycle;
  
  % 3. 绘制按比例缩放的结果图 (深红色加粗)
  % (1) 缩小后的正方形 (1:3 比例 -> 边长 2x2)
  \draw[thick, red!70!black, line width=1.2pt] (10, 6) rectangle (12, 8);
  
  % (2) 放大后的三角形 (2:1 比例 -> 宽4, 高6)
  \draw[thick, red!70!black, line width=1.2pt] (20, 2) -- (20, 8) -- (24, 8) -- cycle;

\end{tikzpicture}
\end{center}

\vspace{0.8cm}

{\small\color{gray}
\textbf{解析：}\\
1. \textbf{图形的缩小}：按 $1:3$ 的比缩小图形，就是把原图形的各边长缩小到原来的 $\frac{1}{3}$。原正方形的边长是 6 格，缩小后得到的新正方形边长应为 $6 \div 3 = 2$ 格。\\
2. \textbf{图形的放大}：按 $2:1$ 的比放大图形，就是把原图形的各边长放大到原来的 2 倍。原直角三角形的竖直直角边长为 3 格，水平直角边长为 2 格。放大后，新的竖直直角边长应为 $3 \times 2 = 6$ 格，新的水平直角边长应为 $2 \times 2 = 4$ 格。最后连接两端点画出新的斜边即可。
}

\vspace{0.6cm}

{\small\color{blue!70!black}
\textbf{绘图基本原理与步骤：}\\
\textbf{1. 统一顶端水平基准线}：经过对参考解答配图的解构，无论是缩小前后的两个正方形，还是放大前后的两个三角形，它们的最高处全部对齐在同一条水平网格线上。因此代码中强制让所有几何图形的高位顶点横向标高均为 \texttt{y=8}（距底边 \texttt{y=0} 共计 \texttt{8} 个单位），以此锚定全图排版的视觉平衡点。\\
\textbf{2. 尺寸方程严格校验}：
   \begin{itemize}
       \item 缩小矩阵：原框 $(2,2) \to (8,8)$，边长跨度为 $6$ 单元。实施算子拉伸 $\times \frac{1}{3}$，生成答案边长跨度为 $2$ 单元。采用 \texttt{rectangle} 指令从 $(10,6)$ 平行绘制至 $(12,8)$。
       \item 放大矩阵：原三角形直角边长为基向量 $(dx=2, dy=3)$。实施放大因子 $\times 2$，其结果图的直角边长强制满足 $(dx=4, dy=6)$。使用全闭合坐标串 \texttt{(20,8) -- (20,2) -- (24,8) -- cycle} 精确构建镜像直角边。
   \end{itemize}
\textbf{3. 专业答卷视觉风格}：所有原题干已知图形统一剥离为黑色常规加粗线型（\texttt{[thick, black]}），所有要求手工绘出的作答结果图层统一覆写辅导读物经典的醒目厚重深红色（\texttt{[thick, red!70!black, line width=1.2pt]}），主次逻辑瞬间明确。
}

\vspace{0.8cm}
\vspace{0.8cm}



\section{解答：只含一条对称轴的双圆组合}

\vspace{0.4cm}

\noindent\textbf{4. 画两个圆，使这两个圆组成的图形只有一条对称轴。}

\vspace{1.2cm}

\begin{center}
% =====================================================================
% 【整体绘制思路：非全对称双圆系统】
% 1. 圆的轴对称性质：圆有无数条对称轴。两个半径相同的圆组合，至少有两条对称轴（连心线及其垂直平分线）。
% 2. 破题关键：只有当两个圆的**半径不同**，且它们**不完全同心**时，由它们组成的组合图形才会丧失中垂线对称性，从而仅保留唯一的一条对称轴——即它们圆心所在的直线。
% 3. 几何构建：
%    - 上方绘制大圆，圆心 O_1 在 (0, 1.5)，半径 1.5。
%    - 下方绘制小圆，圆心 O_2 在 (0, -1)，半径 1。它们是外切关系，相切于 (0, 0)。
%    - 这两个圆组成的“不倒翁”形/雪人形，只有竖直方向(x=0)这一条对称轴。
% =====================================================================
\begin{tikzpicture}[>=Stealth, scale=0.9]
  
  % 1. 绘制核心解域图形 (大圆 + 小圆，采用红色高亮答案线型)
  \draw[thick, red!70!black, line width=1.2pt] (0, 1.5) circle (1.5);
  \draw[thick, red!70!black, line width=1.2pt] (0, -1) circle (1);
  
  % 2. 标记圆心
  \filldraw[black] (0, 1.5) circle (1.5pt) node[right, font=\small, gray] {$O_1$};
  \filldraw[black] (0, -1) circle (1.5pt) node[right, font=\small, gray] {$O_2$};
  
  % 3. 绘制题目要求的本质属性：这唯一的“一条对称轴”
  % 使用蓝色虚线向两端无限延伸
  \draw[dash dot, blue, thick] (0, -2.8) -- (0, 3.8);
  
\end{tikzpicture}
\end{center}

\vspace{0.8cm}

{\small\color{gray}
\textbf{解析：}\\
如果两个圆的大小（半径）完全相同，那么除了穿过两个圆心的那条直线以外，连接两个圆心线段的垂直平分线也是它们的对称轴（此时共有两条互相垂直的对称轴）。\\
题目要求组成的图形“\textbf{只有一条对称轴}”，诀窍就是打破另一种方向的平衡对称。因此我们只能画两个\textbf{大小不同（半径不相等）}的圆。此时，这两个有着一大一小形状的圆所在的整个图形组合，就只能沿它们\textbf{圆心所在的直线}折叠才能全等重合。所以它有且仅有一条对称轴。\\
（注：两个圆可以相离、相切或相交，只要半径不等、圆心不在一点，就都符合题意。）
}

\vspace{0.6cm}

{\small\color{blue!70!black}
\textbf{绘图基本原理与步骤：}\\
\textbf{1. 几何破题设计}：利用 TikZ 函数组建非对称关联域。设置上端大圆圆心锚点为 $(0, 1.5)$ 并赋予半径参量 $r_1=1.5$；下端小圆圆心配置于 $(0, -1)$ 并给定半径 $r_2=1$。大小半径的鲜明对比彻底消灭了图形在水平方向上的镜面反射能力，构建起只存在单一纵向镜像体系的“雪人结构”。\\
\textbf{2. 辅助圆心对齐法}：为了满足“图形形成”这一概念的紧密逻辑，刻意使两个圆精确外切于坐标原点 $(0,0)$。这两个操作生成的数学圆域全部利用辅导教材专供解答高亮色（\texttt{[thick, red!70!black, line width=1.2pt]}）进行粗重轮廓描边。\\
\textbf{3. 核心轴线高亮解构}：图中最关键的题目诉求要素——那条唯一的“对称轴”，则强行剥离图外，采用数学经典标准辅助线型（\texttt{[dash dot, blue, thick]}）垂直贯穿 $O_1$ 与 $O_2$ 连心线并延伸出象限外。这把核心隐藏条件彻底视觉显性化，可一击洞穿本题的内核考点。
}

\vspace{0.8cm}
\vspace{0.8cm}



\section{解答：用数对表示位置}

\vspace{0.4cm}

\noindent\textbf{1. 下面是青莲乡平面图的一部分。}\\
\textbf{(1) 在图上分别用数对表示乡政府、荷花村、海棠村、李庄村和北郭村的位置。}

\vspace{1.2cm}

\begin{center}
% =====================================================================
% 【整体绘制思路：坐标系地图定位】
% 1. 标准直角网格：宽21、高7。全部画为细实线 [thin, black]。
% 2. 坐标数值环绕：在底边和左边标注数字 0-21 和 1-7。
% 3. 定位锚框：绘制右上角 1千米 刻度标尺，以及 N 指北针。
% 4. 村落信息及作答：
%    - 黑点位于给定数对坐标。
%    - 标签文字严格按照原题给定的分布位置和折行排版。
%    - 括弧中的作答数据 (x, y) 采用蓝色字体加浅蓝底色框(\colorbox)渲染，模拟真实高亮作答涂写反馈。
% =====================================================================
\begin{tikzpicture}[>=Stealth, scale=0.6, font=\small]
  % 1. 绘制实线网格
  \draw[step=1cm, black, thin] (0,0) grid (21,7);
  
  % 2. 绘制坐标轴数字
  \foreach \x in {0,1,...,21} {
    \node[below] at (\x, 0) {\x};
  }
  \foreach \y in {1,2,...,7} {
    \node[left] at (0, \y) {\y};
  }
  
  % 3. 右上角刻度尺 (1千米)
  % 使用手绘连接线框代替 package 依赖的大括号，保证任何编译环境安全
  \draw[black] (20, 7.1) -- (20, 7.2) -- (21, 7.2) -- (21, 7.1);
  \node[above] at (20.5, 7.2) {1千米};
  \draw[black] (21.1, 7) -- (21.2, 7) -- (21.2, 6) -- (21.1, 6);
  \node[right] at (21.2, 6.5) {1千米};
  
  % 4. 指北针
  \draw[->, thick] (22.5, 7) -- (22.5, 8) node[above] {N};
  
  % 5. 绘制村庄分布与作答高亮
  % 定义点样式
  \tikzstyle{dot}=[circle, fill=black, inner sep=1.2pt]
  
  % 李庄村 (2, 5)
  \node[dot] at (2, 5) {};
  \node[anchor=west] at (1.2, 4.4) {李庄村({\textcolor{blue}{2, 5}})};
  
  % 乡政府 (3, 1)
  \node[dot] at (3, 1) {};
  \node[anchor=west] at (3.2, 1) {乡政府({\textcolor{blue}{3, 1}})};
  
  % 北郭村 (10, 6)
  \node[dot] at (10, 6) {};
  \node[anchor=north west] at (10, 5.8) {北郭村({\textcolor{blue}{10, 6}})};
  
  % 荷花村 (11, 1)
  \node[dot] at (11, 1) {};
  \node[anchor=west] at (11.2, 1) {荷花村({\textcolor{blue}{11, 1}})};
  
  % 海棠村 (18, 3)
  \node[dot] at (18, 3) {};
  \node[align=center] at (19.5, 3) {海棠村\\[-0.5ex]({\textcolor{blue}{18, 3}})};

\end{tikzpicture}
\end{center}

\vspace{0.8cm}

{\small\color{gray}
\textbf{解析：}\\
数对表示位置的方法是：先表示列数（水平横向的 $x$ 轴坐标），后表示行数（垂直纵向的 $y$ 轴坐标），两个数字中间用逗号隔开，并用小括号括起来。\\
1. \textbf{李庄村}：从左向右数在第 2 列，从下向上数在第 5 行，数对是 $(2, 5)$。\\
2. \textbf{乡政府}：在第 3 列，第 1 行，数对是 $(3, 1)$。\\
3. \textbf{北郭村}：在第 10 列，第 6 行，数对是 $(10, 6)$。\\
4. \textbf{荷花村}：在第 11 列，第 1 行，数对是 $(11, 1)$。\\
5. \textbf{海棠村}：在第 18 列，第 3 行，数对是 $(18, 3)$。
}

\vspace{0.6cm}

{\small\color{blue!70!black}
\textbf{绘图基本原理与步骤：}\\
\textbf{1. 精确构建骨架}：利用 TikZ 中的 \texttt{grid} 指令绘制 $21 \times 7$ 细实线（\texttt{thin, black}）底向网格；遍历 $x$ 轴底端与 $y$ 轴左端分别自动打印步进为 1 的坐标环标；右上角人工利用 \texttt{path} 连线编织防库依赖的阶梯状坐标外置比例尺（标识 $1$ 千米网格尺度），并在右侧外廓配置标准的正北 \texttt{N} 矢量方向箭头。\\
\textbf{2. 坐标位散点刻画}：全部的 5 个村庄点位通过 \texttt{circle, fill=black} 标准圆点标记（锚定极小半径 \texttt{1.2pt}），精准吸附在给定的网格交叉直角坐标映射点上（如 $(2,5)、(10,6)$ 等）。\\
\textbf{3. 作答视觉模拟呈现}：为了最高程度还原教辅资料“印刷体人工填空”的高质感版式，所有村庄文字的 \texttt{node} 节点完全根据参考题干图的视觉重心点进行了单独的偏移量手工微调（利用 \texttt{anchor} 与微距坐标组合），尤其是海棠村实现了精准的跨行居中；填入小括号中的作答核心数据（数字及逗号）采用深蓝色（\texttt{\char`\\textcolor\{blue\}}）字体突出，同时首次引入底色框包边渲染语法（\texttt{\char`\\colorbox\{blue!5\}}）在蓝色文本的背面打上了淡淡的高亮蓝印泥底色，准确地复刻出了原书批阅时的彩色高保真涂写层次效果。
}

\vspace{0.8cm}
\vspace{0.8cm}



\section{解答：方位与比例尺的综合应用}

\vspace{0.4cm}

\noindent\textbf{2. 下面是青山旅游景区平面图的一部分。以香山阁为观测点，填写表格。}

\vspace{0.6cm}

\begin{center}
% =====================================================================
% 【整体绘制思路：极坐标定位与表格综合】
% 1. 坐标系基准：以“香山阁”为原点 (0,0)，建立带有指北针的十字直角坐标系。
% 2. 极坐标映射：
%    - 凤山寺：北偏西 75度 => 从正北(+y)向西逆时针旋转 75度 => 标准极角 90+75 = 165度。图上距离 3cm => (165:3)。
%    - 黄龙洞：南偏东 50度 => 从正南(-y)向东逆时针旋转 50度 => 标准极角 -90+50 = -40度。图上距离 2cm => (-40:2)。
% 3. 装饰点：香山阁、凤山寺、黄龙洞的位置均使用双层同心圆结构（外圈空心，内心实点）标示，精准还原教辅印刷风格。
% 4. 下方铺设 \tabular 表格完成作答填空。
% =====================================================================
\begin{tikzpicture}[>=Stealth, scale=1.2, font=\small]
  
  % 十字基准线与指北针
  \draw[thick] (-3.5, 0) -- (4, 0);
  \draw[->, thick] (0, -2) -- (0, 2) node[above] {N};
  
  % 原点：香山阁 (双重圆圈)
  \filldraw[fill=white, draw=black, thick] (0,0) circle (1.5pt);
  \filldraw[black] (0,0) circle (0.5pt);
  \node[above right, inner sep=3pt] at (0,0) {香山阁};
  
  % 凤山寺 (165度, 距离 3)
  \draw[thick] (0,0) -- (165:3);
  \filldraw[fill=white, draw=black, thick] (165:3) circle (1.5pt);
  \filldraw[black] (165:3) circle (0.5pt);
  \node[above left, inner sep=2pt] at (165:3) {凤山寺};
  
  % 黄龙洞 (-40度, 距离 2)
  \draw[thick] (0,0) -- (-40:2);
  \filldraw[fill=white, draw=black, thick] (-40:2) circle (1.5pt);
  \filldraw[black] (-40:2) circle (0.5pt);
  \node[above right, inner sep=3pt] at (-40:2) {黄龙洞};
  
  % 比例尺文本
  \node[right] at (2, -1.5) {比例尺 $1:20000$};

\end{tikzpicture}
\end{center}

\vspace{0.4cm}

% 填空表格（根据要求：只输出蓝色字体作答，不加背景阴影）
\begin{center}
\renewcommand{\arraystretch}{1.6}
\begin{tabular}{c|c|c|c}
\noalign{\hrule height 1pt}
\makebox[2cm][c]{景\quad 点} & \makebox[4.5cm][c]{方\quad 向} & \makebox[2.5cm][c]{图上距离/cm} & \makebox[2.5cm][c]{实际距离/m} \\
\hline
凤山寺 & (\ \textcolor{blue}{北}\ )\,偏\,(\ \textcolor{blue}{西}\ )\,(\ \textcolor{blue}{75}\ )$^{\circ}$ & \textcolor{blue}{3} & \textcolor{blue}{600} \\
\hline
黄龙洞 & (\ \textcolor{blue}{南}\ )\,偏\,(\ \textcolor{blue}{东}\ )\,(\ \textcolor{blue}{50}\ )$^{\circ}$ & \textcolor{blue}{2} & \textcolor{blue}{400} \\
\noalign{\hrule height 1pt}
\end{tabular}
\end{center}

\vspace{0.8cm}

{\small\color{gray}
\textbf{解析：}\\
1. \textbf{确定方向}：以香山阁为观测点。凤山寺在香山阁的左上方，也就是北偏西方向，根据量角器测量或题目隐性设定，角度为 $75^{\circ}$；黄龙洞在香山阁的右下方，也就是南偏东方向，角度为 $50^{\circ}$。\\
2. \textbf{测量并计算距离}：测量图上距离，凤山寺到香山阁的图上距离是 $3\text{ cm}$，黄龙洞到香山阁的图上距离是 $2\text{ cm}$。\\
根据比例尺 $1:20000$，即图上 $1\text{ cm}$ 代表实际距离 $20000\text{ cm} = 200\text{ m}$。\\
所以，凤山寺的实际距离 $= 3 \times 200 = 600\text{ (m)}$；\\
黄龙洞的实际距离 $= 2 \times 200 = 400\text{ (m)}$。
}

\vspace{0.6cm}

{\small\color{blue!70!black}
\textbf{绘图与制表基本原理与步骤：}\\
\textbf{1. 极坐标空间建模}：以观测点构筑局部直角坐标定位十字准线。对于已知方向和距离的点位延伸，使用极坐标 \texttt{(angle:radius)} 是最高效的数学定位法。凤山寺的“北偏西 $75^{\circ}$”等价于直角坐标系的极角 $90^{\circ}+75^{\circ}=165^{\circ}$，配合图上量取的 $3\text{cm}$ 长度直接绘制向量 \texttt{(165:3)}；黄龙洞的“南偏东 $50^{\circ}$”等价于 $-90^{\circ}+50^{\circ}=-40^{\circ}$，配合 $2\text{cm}$ 长度构成向量 \texttt{(-40:2)}。\\
\textbf{2. 标点符号细节复刻}：所有定位点均采用了“外圆内点”（双重节点覆盖图案）的设计（\texttt{fill=white, draw=black} 在圆体内嵌套微小实心 \texttt{black} 黑点），从而精准还原考卷印刷时地理测量指示图表质感。\\
\textbf{3. 作答数据表格极简渲染}：下挂的作答表格通过标准 \texttt{tabular} 环境全边框构建，通过 \texttt{\char`\\makebox} 严格控宽来排版表格表头（如“景\quad 点”），规避了默认挤压。特别遵照“制表，但是不用加蓝色阴影”的明确需求指令，摒弃了 \texttt{\char`\\colorbox} 底纹函数，仅使用 \texttt{\char`\\textcolor\{blue\}} 将提取出的计算结果以原笔迹高亮蓝色空灵嵌入括号留白内，完全匹配目标要求。
}

\vspace{0.8cm}


\vspace{0.6cm}

{\small\color{blue!70!black}
\textbf{绘图基本原理与步骤：}\\
\textbf{1. 坐标定位与几何构建}：根据题意中的已知长度和角度，使用 \texttt{coordinate} 定义关键顶点坐标，通过 \texttt{draw} 指令按序连接成完整几何图形。线宽、颜色参数区分原图（黑色）与答案（红色 \texttt{red!70!black}）图层。\\
\textbf{2. 辅助量标注}：利用 \texttt{node} 的方向锚点（\texttt{above}、\texttt{below}、\texttt{left}、\texttt{right}）将角度、边长、面积等数据标注在图形周围，避免与线条或其他标注重叠。若涉及角弧则使用 \texttt{arc} 或 \texttt{pic[angle]} 命令渲染。\\
\textbf{3. 虚实线分层}：题干已知条件用实线绘制，解答新增的辅助线或操作痕迹用 \texttt{dashed} 虚线或彩色加粗线标识，确保读者一眼区分原有信息与作答内容。
}

%% ==============================
%% 六年级统计表：1分钟跳绳成绩分段统计
%% ==============================
\newpage




\section{解答：图形的放大与缩小}

\vspace{0.4cm}

\noindent\textbf{5. 按 2:1 的比画出长方形放大后的图形，再按 1:3 的比画出三角形缩小后的图形。}

\vspace{0.8cm}

% =====================================================================
% 【整体绘制思路：图形的放大与缩小】
% 1. 坐标系与网格环境：
%    使用 step=1.0 的虚线在 20x9 的区间绘制方格纸背景，左下角定为坐标原点 (0,0)。
% 2. 比例计算体系：
%    - 原长方形：长 3、宽 1。放大比 2:1 -> 新图长 6、宽 2。
%    - 原三角形：底 9、高 6。缩小比 1:3 -> 新图底 3、高 2。
% 3. 定位布局（复刻参考答案视觉）：
%    - 原长方形置于左上方，新长方形（红框）置于左下方。
%    - 原三角形置于中右方，新三角形（红框）置于其右上方空白处。
% 4. 色彩映射逻辑：
%    - 原题干上的已有图形使用 thick, black 绘制。
%    - 学生需要作答的线条使用 thick, red 绘制，凸显“作图题答案”。
% =====================================================================

\begin{center}
\begin{tikzpicture}[>=Stealth, line width=1.0pt, scale=0.8]

  % =========== 背景网格 ===========
  % \draw[step=1.0, ...]：每隔 1 个单位画一条线
  % grid (20,9)：从 (0,0) 画到 (20,9)，横向 20 个格子，纵向 9 个格子
  % gray!70, dashed：灰色虚线，模拟常见的方格纸效果
  \draw[step=1.0, gray!70, dashed, very thin] (0,0) grid (20,9);
  % 画网格最外侧的实线边框
  \draw[gray!80, thin] (0,0) rectangle (20,9);

  % =========== 原图形（题干黑色图形） ===========
  % 1. 原始长方形 (长 3, 宽 1)
  % 放在左上角，坐标从 (1,7) 到 (4,8)
  \draw[black] (1,7) rectangle (4,8);

  % 2. 原始直角三角形 (底 9, 高 6)
  % 左侧直角边对齐 x=9，顶部对齐 y=8 (与长方形顶端齐平)
  \coordinate (T_top) at (9,8);     % 顶点，位于上方
  \coordinate (T_bl)  at (9,2);     % 直角顶点，位于下方
  \coordinate (T_br)  at (18,2);    % 锐角顶点，位于右侧 (9+9=18)
  % cycle 自动连接 (18,2) 并闭合至起点 (9,8)
  \draw[black] (T_top) -- (T_bl) -- (T_br) -- cycle;


  % =========== 作答图形（红色缩放图形） ===========
  % 3. 放大后的长方形：按 2:1 放大 (长 6, 宽 2)
  % 放置在原长方形下方（间隔 1 格），坐标从 (1,4) 到 (7,6)
  \draw[red] (1,4) rectangle (7,6);

  % 4. 缩小后的三角形：按 1:3 缩小 (底 3, 高 2)
  % 放置在原三角形的右上方，顶部平起，即顶部 y=8
  \coordinate (T_ans_top) at (14,8);    % 顶点，y=8
  \coordinate (T_ans_bl)  at (14,6);    % 直角顶点，高 2，所以 y=8-2=6
  \coordinate (T_ans_br)  at (17,6);    % 锐角顶点，底 3，所以右侧 = 14+3=17
  \draw[red] (T_ans_top) -- (T_ans_bl) -- (T_ans_br) -- cycle;

\end{tikzpicture}
\end{center}

\vspace{0.6cm}

{\small\color{gray}
\textbf{解析：}\\
1. **长方形的放大**：测量原长方形可知其长为 $3$ 个方格，宽为 $1$ 个方格。按照 $2:1$ 的比放大，也就是将边长扩大到原来的 $2$ 倍。放大后的长为 $3 \times 2 = 6$ 个方格，宽为 $1 \times 2 = 2$ 个方格。据此在大网格中画出新的 $6 \times 2$ 长方形。\\
2. **三角形的缩小**：测量原直角三角形可知其两条直角边（底和高）分别为 $9$ 个方格和 $6$ 个方格。按照 $1:3$ 的比缩小，即边长缩小为原来的 $\frac{1}{3}$。缩小后的底为 $9 \div 3 = 3$ 个方格，高为 $6 \div 3 = 2$ 个方格。选取网格右侧空白处，以 $3$、 $2$ 为直角边画出一个符合原图姿态的直角三角形。
}

\vspace{0.6cm}

{\small\color{blue!70!black}
\textbf{绘图基本原理与步骤：}\\
\textbf{1. 虚线网格系统搭建}：运用 \texttt{grid} 命令配合 \texttt{dashed} 属性，构建底层的 $20 \times 9$ 网格，为线段测量和比例计算提供像素级准确的参照系。\\
\textbf{2. 解构题设与变量转换}：首先固定题干已给图形（\texttt{black} 实线），获取原尺寸（$3\times 1$ 与 $9\times 6$）；随后分别进行代数倍乘与除法（放大倍数与缩小比例），得出目标尺寸（$6\times 2$ 与 $3\times 2$）。\\
\textbf{3. 多元素布局防重叠}：通过对全局画布（$0 \le x \le 20, 0 \le y \le 9$）进行预规划：将大面积答案（大红框）沉于左侧中部，小面积答案（小红三角）悬于右上，既满足紧凑要求亦实现与参考答案在视觉落点上的高度一致。\\
\textbf{4. 定向颜色标注}：严格施行“原题黑、答案红”的双色区分法，强化辅助教材的直观导向性。
}

\newpage
